\chapter{Modul atas Daerah Ideal Utama}\label{ch:b3:modules}

Aljabar linear atas sebuah \emph{ring} alih-alih atas sebuah lapangan: perubahan
hipotesis yang kecil ini melahirkan salah satu teorema penyatuan
terbesar dalam aljabar. \omterm{def:b3:modules:module}{Modul} atas $\Z$ adalah grup abelian; \omterm{def:b3:modules:module}{modul} atas
$K[X]$ adalah ruang vektor yang dilengkapi sebuah endomorfisma. Karena itu
teorema struktur untuk \omterm{def:b3:modules:module}{modul} yang \omterm{def:b3:modules:free}{dibangun secara berhingga} atas sebuah \omterm{def:b3:rings:pidufd}{DIU}
mengklasifikasikan, dalam satu tarikan, semua grup abelian yang \omterm{def:b3:modules:free}{dibangun secara
berhingga} \emph{sekaligus} semua endomorfisma sampai pada keserupaan ---
reduksi Jordan, yang pada Tahun ke-2 diperoleh lewat induksi yang berliku,
jatuh sebagai akibat, bersama saudaranya yang lebih halus, yaitu
bentuk kanonik \emph{rasional} yang berlaku atas setiap lapangan. Mesin
perhitungannya adalah \omterm{thm:b3:modules:smith}{bentuk normal Smith}, sebuah aritmetika
matriks yang layak disandingkan dengan Euclid.

Sepanjang bab ini $A$ adalah ring komutatif, yang segera menjadi \omterm{def:b3:rings:pidufd}{DIU}; kata ``\omterm{def:b3:modules:module}{modul}''
berarti modul-$A$.

\section{Modul dan modul bebas}

\begin{definition}\label{def:b3:modules:module}
Sebuah \emph{modul-$A$}\index{modul} adalah grup abelian $(M, +)$
dengan perkalian skalar $A \times M \to M$ yang memenuhi
aksioma ruang vektor: $a(x + y) = ax + ay$, $(a + b)x = ax + bx$,
$(ab)x = a(bx)$, $1x = x$. Submodul, modul kuosien $M/N$, morfisma
(pemetaan linear-$A$), jumlah langsung $\bigoplus_i M_i$, dan
teorema isomorfisma didefinisikan dan dibuktikan kata demi kata seperti pada
ruang vektor dan grup abelian; khususnya $M/\ker f \cong
\operatorname{im} f$ untuk sebuah morfisma $f$.
\end{definition}

\begin{example}\label{ex:b3:modules:examples}
Berikut tiga kasus yang memotivasinya.
\begin{enumerate}
  \item $A = K$ sebuah lapangan: \omterm{def:b3:modules:module}{modulnya} adalah ruang vektor.
  \item $A = \Z$: \omterm{def:b3:modules:module}{modulnya} tepat berupa grup abelian (sebab $nx$
        terpaksa sama dengan $x + \dots + x$), dan submodulnya adalah subgrup.
  \item $A = K[X]$: \omterm{def:b3:modules:module}{modulnya} berupa ruang vektor-$K$ $V$ beserta
        pemetaan linear-$K$ $u\colon x \mapsto X\cdot x$ ---
        sebaliknya, setiap pasangan $(V, u)$ dengan $u \in \mathcal L(V)$
        menjadi modul-$K[X]$ lewat $P \cdot x = P(u)(x)$.
        Submodulnya tepat berupa subruang yang stabil terhadap $u$.
\end{enumerate}
\omterm{def:b3:rings:ideal}{Ideal} $A$ tidak lain adalah submodul $A$; \omterm{def:b3:rings:ideal}{ring kuosien}
$A/I$ merupakan modul-$A$. Berbeda dengan ruang vektor, \omterm{def:b3:modules:module}{modul} dapat mempunyai
\emph{\omterm{def:b3:modules:torsion}{torsi}}: di $\Z/6\Z$, unsur $\bar 3 \ne 0$ dimatikan
oleh $2 \neq 0$.
\end{example}

\begin{definition}\label{def:b3:modules:free}
\omterm{def:b3:modules:module}{Modul} $M$ disebut \emph{dibangun secara berhingga}\index{modul dibangun secara berhingga} jika
$M = Ax_1 + \dots + Ax_n$ untuk suatu $x_i$. \omterm{def:b3:modules:module}{Modul} $M$ disebut
\emph{bebas}\index{modul bebas} berrank $n$ jika $M \cong A^n$,
yakni jika ia mempunyai \emph{basis} (keluarga pembangun yang
bebas linear atas $A$). Setiap $M$ yang dibangun secara berhingga merupakan
kuosien sebuah \omterm{def:b3:modules:module}{modul} bebas: $(a_1, \dots, a_n) \mapsto \sum a_ix_i$
memetakan $A^n$ pada $M$.
\end{definition}

\begin{proposition}[Kekekalan rank]\label{prop:b3:modules:rank}
Jika $A \neq 0$ dan $A^m \cong A^n$, maka $m = n$.
\end{proposition}

\begin{proof}
Pilih sebuah \omterm{def:b3:rings:primemaximal}{ideal maksimal} $\mathfrak m$ dari $A$ (\cref{thm:b3:rings:krull})
lalu tetapkan $k = A/\mathfrak m$, sebuah lapangan. Isomorfisma $f \colon A^m
\to A^n$ memetakan $\mathfrak m A^m$ ke dalam $\mathfrak m A^n$
(karena linearitas), sehingga menginduksikan isomorfisma antar-kuosien
\[
A^m/\mathfrak m A^m \;\cong\; A^n/\mathfrak m A^n,
\qquad\text{yakni}\qquad k^m \cong k^n
\]
sebagai ruang vektor-$k$ (kuosien $A^m/\mathfrak m A^m$
dimatikan oleh $\mathfrak m$, sehingga aksi $A$ terfaktorkan melalui $k$;
peta basis bakunya membentuk basis-$k$). Teori dimensi
atas lapangan $k$ lalu memberikan $m = n$.
\end{proof}

\begin{theorem}[Submodul dari modul bebas]\label{thm:b3:modules:submodule}
Misalkan $A$ sebuah \omterm{def:b3:rings:pidufd}{DIU} dan $M \subseteq A^n$ sebuah submodul. Maka $M$
\omterm{def:b3:modules:free}{bebas} dengan rank $\leq n$.
\end{theorem}

\begin{proof}
Induksi pada $n$. Untuk $n = 1$: $M$ berupa \omterm{def:b3:rings:ideal}{ideal}, jadi $M = (0)$
(\omterm{def:b3:modules:free}{bebas} berrank $0$) atau $M = dA \cong A$ (sebab $x \mapsto dx$
injektif di daerah integral). Untuk $n > 1$: misalkan $\pi \colon A^n \to A$
koordinat terakhirnya. Maka $\pi(M)$ berupa \omterm{def:b3:rings:ideal}{ideal}, entah $(0)$ entah $dA$. Jika
$(0)$: $M \subseteq A^{n-1} \times \{0\}$ dan induksinya berlaku.
Jika tidak, pilihlah $x_0 \in M$ dengan $\pi(x_0) = d$. Setiap $x \in M$
ditulis secara tunggal sebagai
\[
x = \underbrace{\frac{\pi(x)}{d}}_{\in A}\, x_0 +
\Bigl(x - \tfrac{\pi(x)}d x_0\Bigr),
\qquad x - \tfrac{\pi(x)}d x_0 \in M \cap \ker\pi
\]
(sebab $\pi(x) \in dA$, sehingga koefisiennya di $A$). Dengan demikian $M = Ax_0
\oplus (M \cap \ker \pi)$: jumlahnya langsung sebab $\pi(ax_0) =
ad = 0$ memaksa $a = 0$. Menurut induksi, $M \cap \ker\pi \subseteq
\ker \pi \cong A^{n-1}$ \omterm{def:b3:modules:free}{bebas} dengan rank $\leq n - 1$; menambahkan
$x_0$ (yang \omterm{def:b3:modules:free}{bebas} dari $\ker\pi$ seperti baru saja terlihat) memberikan basis
$M$ yang berkardinalitas $\leq n$.
\end{proof}

\begin{remark}\label{rem:b3:modules:presentation}
Akibatnya, atas sebuah \omterm{def:b3:rings:pidufd}{DIU} setiap \omterm{def:b3:modules:module}{modul} $M$ yang \omterm{def:b3:modules:free}{dibangun secara berhingga} mempunyai
\emph{penyajian berhingga}: surjeksi $\varphi\colon A^n \to M$
mempunyai kernel \omterm{def:b3:modules:free}{bebas} dengan basis $c_1, \dots, c_k$ ($k \leq n$), dan $M
\cong A^n / \operatorname{im}(C)$ dengan $C \in M_{n,k}(A)$ berupa
matriks yang kolomnya adalah $c_j$. Memahami $M$ berarti
memahami sebuah \emph{matriks atas $A$} sampai pada pergantian basis di
sumber dan sasarannya --- itulah pokok bagian berikutnya.
\end{remark}

\section{Bentuk normal Smith}

\begin{definition}\label{def:b3:modules:equivalence}
Dua matriks $B, C \in M_{n,k}(A)$ disebut \emph{ekuivalen} jika $C =
QBP$ dengan $Q \in GL_n(A)$, $P \in GL_k(A)$ (yang terbalikkan
\emph{atas $A$}: determinannya di $A^\times$). Matriks penyajian
yang ekuivalen mendefinisikan \omterm{def:b3:modules:module}{modul} $A^n/
\operatorname{im}$ yang isomorfik (ganti basis di $A^n$ dan $A^k$).
\end{definition}

\begin{theorem}[Bentuk normal Smith]\label{thm:b3:modules:smith}
Misalkan $A$ sebuah \omterm{def:b3:rings:pidufd}{DIU} dan $B \in M_{n,k}(A)$. Maka $B$ ekuivalen dengan
sebuah matriks diagonal\index{bentuk normal Smith}
\[
\operatorname{diag}(d_1, d_2, \dots, d_r, 0, \dots, 0),
\qquad d_1 \mid d_2 \mid \cdots \mid d_r \neq 0 ,
\]
dan $d_i$ tunggal sampai pada kesekawanan: $d_1 \cdots d_i$ merupakan
FPB semua minor $i \times i$ dari $B$ (khususnya FPB itu
merupakan invarian keekuivalenan). Bilangan $d_i$ disebut \emph{faktor
invarian}\index{faktor invarian} dari $B$.
\end{theorem}

\begin{proof}
\emph{Keberadaan.} Jika $B = 0$, selesai. Jika tidak, tinjau himpunan
\omterm{def:b3:rings:ideal}{ideal} $(b)$ yang dibangun oleh entri matriks yang ekuivalen dengan $B$;
karena $A$ bersifat \omterm{def:b3:rings:noetherian}{Noether}, pilihlah matriks $B'$ yang ekuivalen dengan $B$
dan sebuah entri $d$ dari $B'$ dengan $(d)$ \emph{maksimal} di himpunan itu.
Pindahkan $d$ ke posisi $(1,1)$ lewat penukaran baris dan kolom.

\emph{Klaim: $d$ membagi setiap entri $B'$.} Pertama, kolom 1: jika
$b_{i1}$ bukan kelipatan $d$, misalkan $e = \gcd(d, b_{i1}) = ud +
vb_{i1}$ (Bézout), sehingga $(e) \supsetneq (d)$. Muslihat matriks
$2 \times 2$: bekerja pada baris $1$ dan $i$ dengan
\[
\begin{pmatrix} u & v \\ -\dfrac{b_{i1}}e & \dfrac de
\end{pmatrix}
\in GL_2(A)
\qquad
\Bigl(\det = \tfrac{ud + v b_{i1}}e = 1\Bigr)
\]
menghasilkan matriks ekuivalen dengan entri $e$ pada posisi $(1,1)$:
dan ini bertentangan dengan kemaksimalan $(d)$. Jadi $d$ membagi kolom $1$,
dan secara simetris juga baris $1$. Mengurangkan kelipatan baris 1 dan
kolom 1 membersihkan keduanya: jadi $B'$ ekuivalen dengan $\begin{pmatrix} d &
0\\ 0 & B''\end{pmatrix}$. Selanjutnya, $d$ membagi setiap entri $b$ dari
$B''$: tambahkan baris tempat $b$ berada pada baris 1 (sebuah operasi elementer; baris
pertama yang baru memuat $d$ dan entri $b$), lalu ulangi
argumen pembersihan kolomnya: entri yang bukan kelipatannya akan kembali memperbaiki
$(d)$. Sekarang berinduksilah pada ukurannya: $B''$, yang semua entrinya
habis dibagi $d$, mempunyai bentuk Smith $\operatorname{diag}(d_2,
\dots)$ yang entrinya tetap habis dibagi $d$ (setiap entri sembarang
$QB''P$ merupakan kombinasi-$A$ dari entri $B''$); tetapkan $d_1 =
d$.

\emph{Ketunggalan.} Misalkan $D_i(B)$ melambangkan FPB semua minor $i \times i$.
Operasi baris dan kolom, dan lebih umum lagi
perkalian dengan \emph{sembarang} matriks, tak dapat memperkecil FPB itu: sebab
minor $i \times i$ dari $QB$ merupakan kombinasi-$A$ dari minor $B$
(penjabaran Cauchy--Binet; atau secara langsung: setiap baris $QB$ merupakan
kombinasi baris $B$, dan minor bersifat multilinear terhadap barisnya).
Jadi $D_i(QBP)$ dan $D_i(B)$ saling membagi: dengan kata lain $D_i$ merupakan
invarian keekuivalenan. Pada bentuk diagonalnya, minor $i
\times i$ yang tak nol berupa hasil kali $i$ buah $d_j$, dan
keterbagian $d_1 \mid \dots \mid d_r$ menjadikan $d_1 \cdots d_i$ sebagai
FPB-nya. Karena itu $d_1 \cdots d_i = D_i(B)$ sampai pada unit, dan $d_i =
D_i/D_{i-1}$ tertentu.
\end{proof}

\begin{method}\label{met:b3:modules:smith}
Atas \omterm{def:b3:rings:pidufd}{daerah Euclid} ($\Z$, $K[X]$), reduksi Smith merupakan sebuah
algoritma --- tanpa perlu argumen kemaksimalan: bawalah entri dengan
ukuran Euclid terkecil ke posisi $(1,1)$; jika ia tak membagi
suatu entri pada baris atau kolomnya, pembagian Euclid meninggalkan
sisa yang lebih kecil di sana --- tukarkan sisa itu ke posisi itu lalu mulai lagi
(prosesnya berhenti karena ukurannya menurun); bila ia sudah membagi seluruh baris dan
kolomnya, bersihkan keduanya; jika ia tak membagi entri di dalam, tambahkan baris
itu pada baris $1$ lalu mulai lagi; kemudian rekursi pada blok dalamnya. Dalam
praktik pada matriks bulat: hitunglah $D_1 = \gcd$ entrinya, lalu
$D_2$, \dots\ lewat minor untuk ukuran kecil, atau jalankan algoritmanya.
\end{method}

\begin{example}[Sebuah reduksi Smith, selengkapnya]
\label{ex:b3:modules:smithworked}
Reduksikan $M = \begin{pmatrix} 1 & 2 & 3\\ 4 & 5 & 6\\ 7 & 8 & 9
\end{pmatrix}$ atas $\Z$. Sudutnya, yaitu $1$, membagi segalanya:
bersihkan baris dan kolomnya ($L_2 \leftarrow L_2 - 4L_1$, $L_3
\leftarrow L_3 - 7L_1$, lalu $C_2 \leftarrow C_2 - 2C_1$, $C_3
\leftarrow C_3 - 3C_1$):
\[
M \sim \begin{pmatrix} 1 & 0 & 0\\ 0 & -3 & -6\\ 0 & -6 & -12
\end{pmatrix} .
\]
Pada blok dalamnya, sudutnya $-3$ membagi semua entri: $L_3
\leftarrow L_3 - 2L_2$ dan $C_3 \leftarrow C_3 - 2C_2$ membersihkannya
menjadi $\operatorname{diag}(-3, 0)$. Setelah tandanya disesuaikan (kalikan sebuah
baris dengan $-1$, sebuah operasi yang sah):
\[
M \sim \operatorname{diag}(1, 3, 0),
\qquad
\Z^3/M\Z^3 \cong \Z/3\Z \times \Z .
\]
Periksa silang lewat pembagi determinan: $D_1 =
\gcd(\text{entri}) = 1$; setiap minor $2\times2$ dari $M$ merupakan
kelipatan $3$ (misalnya $\det\bigl(\begin{smallmatrix}1 & 2\\ 4
& 5\end{smallmatrix}\bigr) = -3$) dan salah satunya sama dengan $-3$: jadi $D_2 =
3$; dan $D_3 = \det M = 0$. Karena itu $d_1 = 1$, $d_2 = 3$, $d_3 = 0$:
jawaban yang sama. Dua pelajarannya: \omterm{thm:b3:modules:smith}{faktor invarian} yang nol merekam
turunnya rank (kokernelnya memperoleh suku \omterm{def:b3:modules:free}{bebas} $\Z$), dan
rantai keterbagian $1 \mid 3 \mid 0$ merupakan sertifikat Smith
--- reduksi diagonal yang melanggar rantai itu (katakanlah
$\operatorname{diag}(2, 3)$, yang bisa dihasilkan oleh yang ceroboh
dari $\bigl(\begin{smallmatrix}2 & 0\\ 0 &
3\end{smallmatrix}\bigr)$ karena berhenti terlalu dini: bentuk Smith yang benar
adalah $\operatorname{diag}(1, 6)$, sebab di sini $D_1 = 1$!) belumlah
selesai.
\end{example}

\begin{omfigure}
\begin{tikzpicture}[scale=0.72]
  \clip (-3.4,-1.3) rectangle (4.4,4.3);
  \foreach \x in {-3,...,4}
    \foreach \y in {-1,...,4}
      \fill[black!40] (\x,\y) circle (1.6pt);
  \fill[omDef!12, opacity=0.7] (0,0) -- (1,3) -- (3,3) -- (2,0)
    -- cycle;
  \foreach \a/\b in {0/0, 2/0, 4/0, -2/0, 1/3, 3/3, -1/3, -3/3,
                     2/6, -2/-6, 0/6}
    \fill[omThm] (\a,\b) circle (2.6pt);
  \draw[omThm, thick, ->] (0,0) -- (2,0)
    node[midway, below, font=\small] {$(2,0)$};
  \draw[omThm, thick, ->] (0,0) -- (1,3)
    node[midway, left, font=\small] {$(1,3)$};
\end{tikzpicture}

{\small Subkisi $L = \Z(2,0) + \Z(1,3)$ dari $\Z^2$ (titik
merah). \omterm{thm:b3:modules:smith}{Bentuk normal Smith} dari $\bigl(\begin{smallmatrix} 2 &
1\\ 0 & 3\end{smallmatrix}\bigr)$ adalah
$\operatorname{diag}(1, 6)$: pada basis tersesuaikan $f_1 = (1,3)$,
$f_2 = (0,1)$ dari $\Z^2$, berlaku $L = \Z f_1 \oplus \Z\,6f_2$, sehingga
$\Z^2/L \cong \Z/6\Z$ --- indeksnya sama dengan $\abs{\det} = 6$, yaitu
luas daerah fundamental yang diarsir.}
\end{omfigure}

\section{Teorema struktur}

\begin{definition}\label{def:b3:modules:torsion}
Misalkan $A$ daerah integral dan $M$ sebuah modul-$A$. Submodul
\emph{torsi}\index{torsi} adalah
\[
T(M) = \{x \in M : ax = 0 \text{ untuk suatu } a \neq 0\}
\]
(sebuah submodul: jika $ax = by = 0$ maka $ab(x + y) = 0$ dengan $ab \ne 0$).
\omterm{def:b3:modules:module}{Modul} $M$ disebut \emph{\omterm{def:b3:modules:free}{bebas} torsi} jika $T(M) = 0$, dan \emph{\omterm{def:b3:modules:module}{modul} torsi}
jika $T(M) = M$.
\end{definition}

\begin{theorem}[Struktur modul yang dibangun secara berhingga atas sebuah
DIU]\label{thm:b3:modules:structure}
Misalkan $A$ sebuah \omterm{def:b3:rings:pidufd}{DIU} dan $M$ modul-$A$ yang \omterm{def:b3:modules:free}{dibangun secara berhingga}. Terdapat
tepat satu $r \in \N$ dan unsur tak nol bukan unit $d_1 \mid d_2 \mid
\cdots \mid d_s$, tunggal sampai pada kesekawanan, dengan\index{teorema
struktur modul}
\[
M \;\cong\; A^r \,\oplus\, A/(d_1) \oplus \cdots \oplus A/(d_s).
\]
Lebih lanjut $T(M) \cong A/(d_1)\oplus\dots\oplus A/(d_s)$ dan $M/T(M)
\cong A^r$: jadi \omterm{def:b3:modules:module}{modul} \emph{\omterm{def:b3:modules:free}{bebas} \omterm{def:b3:modules:torsion}{torsi}} yang \omterm{def:b3:modules:free}{dibangun secara berhingga} atas
\omterm{def:b3:rings:pidufd}{DIU} pastilah \omterm{def:b3:modules:free}{bebas}.
\end{theorem}

\begin{proof}
\emph{Keberadaan.} Sajikan $M \cong A^n/\operatorname{im}(C)$
(\cref{rem:b3:modules:presentation}) lalu bawalah $C$ ke bentuk Smith:
setelah kedua pergantian basis, $M \cong A^n / (d_1A \times \dots
\times d_rA \times 0 \times \dots \times 0) \cong A/(d_1) \oplus
\dots \oplus A/(d_r) \oplus A^{\,n - r}$. Buanglah suku
yang $d_i$-nya berupa unit (sebab $A/(d_i) = 0$); rantai keterbagiannya
tetap bertahan.

\emph{Pengenalan \omterm{def:b3:modules:torsion}{torsinya}.} Pada penguraian itu, $A^r$ bersifat
\omterm{def:b3:modules:free}{bebas} \omterm{def:b3:modules:torsion}{torsi} (daerah integral tak punya pembagi nol) dan setiap $A/(d_i)$
bersifat \omterm{def:b3:modules:torsion}{torsi} (dimatikan oleh $d_i \ne 0$); jumlah langsung memecah \omterm{def:b3:modules:torsion}{torsinya}
sejalan dengan itu: $T(M) = \bigoplus_i A/(d_i)$ dan $M/T(M) \cong A^r$.

\emph{Ketunggalan $r$}: $M/T(M) \cong A^r$ hanya bergantung pada $M$,
dan \cref{prop:b3:modules:rank} memakukan nilai $r$.

\emph{Ketunggalan $d_i$}: cukup ditangani \omterm{def:b3:modules:torsion}{modul torsi}
$T = T(M)$. Uraikan setiap $d_i$ atas \omterm{def:b3:rings:divisibility}{unsur prima} lalu pecahlah lewat
teorema sisa Tiongkok (\cref{thm:b3:rings:crt}; \omterm{def:b3:rings:divisibility}{unsur prima} yang berbeda
membangun \omterm{def:b3:rings:ideal}{ideal} yang \omterm{thm:b3:rings:crt}{komaksimal}):
\[
A/(d) \cong \bigoplus_{p \mid d} A/\bigl(p^{v_p(d)}\bigr):
\qquad
T \cong \bigoplus_{p} \bigoplus_{j} A/\bigl(p^{k_{p,j}}\bigr),
\]
yaitu \emph{pembagi elementer}\index{pembagi elementer}
$p^{k_{p,j}}$. Sebaliknya $d_i$ dapat disusun ulang dari
multihimpunan \omterm{thm:b3:modules:structure}{pembagi elementernya} ($d_s$ = hasil kali pangkat tertinggi
setiap \omterm{def:b3:rings:divisibility}{unsur prima}, dan seterusnya), sehingga cukup dibuktikan bahwa
multihimpunan $\{k_{p,j}\}_j$ ditentukan oleh $T$ untuk setiap \omterm{def:b3:rings:divisibility}{unsur prima}
$p$. Tetapkan $p$; untuk $j \geq 1$ tinjaulah ruang vektor-$A/(p)$
$p^{j-1}T/p^jT$. Pada suku siklik $A/(p^k)$:
\[
p^{j-1}\bigl(A/(p^k)\bigr)\big/p^{j}\bigl(A/(p^k)\bigr)
\cong
\begin{cases}
A/(p) & \text{jika } j \leq k,\\
0 & \text{jika } j > k,
\end{cases}
\]
sedangkan pada suku $A/(q^k)$ dengan $q \neq p$: perkalian dengan $p$ di sana
bijektif ($p$ terbalikkan modulo $q^k$ menurut Bézout), sehingga
kuosiennya $0$. Jumlah langsung meneruskan hal ini: $\dim_{A/(p)}
p^{j-1}T/p^jT = \#\{i : k_{p,i} \geq j\}$. Dimensi intrinsik ini
menentukan multihimpunan eksponennya.
\end{proof}

\begin{corollary}[Grup abelian yang dibangun secara berhingga]
\label{cor:b3:modules:abelian}
Setiap grup abelian yang \omterm{def:b3:modules:free}{dibangun secara berhingga} berbentuk $\Z^r \times
\Z/d_1\Z\times\dots\times\Z/d_s\Z$ dengan $d_1 \mid \dots \mid
d_s$, secara tunggal. Setiap grup abelian \emph{berhingga} merupakan hasil kali
grup siklik yang berorde pangkat prima, tunggal sebagai
multihimpunan.\index{grup abelian, struktur}
\end{corollary}

\begin{example}\label{ex:b3:modules:abelian}
Grup abelian berorde $p^n$ bersesuaian dengan \emph{partisi}
dari $n$: untuk $p^4$ ada $\Z/p^4$, $\Z/p^3\times\Z/p$,
$(\Z/p^2)^2$, $\Z/p^2 \times (\Z/p)^2$, $(\Z/p)^4$ --- lima
grup, sebab $4$ mempunyai lima partisi. Orde campuran mengalikan
cacahannya prima demi prima (lewat sisa Tiongkok): jadi ada $5 \times 2$ grup abelian
berorde $2^4 \cdot 3^2 = 144$.
\end{example}

\section{Terapan: bentuk kanonik endomorfisma}

Misalkan $K$ sebuah lapangan, $V$ ruang vektor-$K$ berdimensi berhingga
$n$, dan $u \in \mathcal L(V)$; jadikan $V$ modul-$K[X]$ lewat $P
\cdot x = P(u)(x)$ (\cref{ex:b3:modules:examples}). \omterm{def:b3:modules:module}{Modul} ini
\omterm{def:b3:modules:free}{dibangun secara berhingga} (basis-$K$ membangunnya) dan bersifat \omterm{def:b3:modules:torsion}{torsi}: untuk
setiap $x$, kesatuan $n+1$ vektor $x, u(x), \dots, u^n(x)$ bergantung
secara linear atas $K$, sehingga memberikan polinomial penghapus yang tak nol.

\begin{definition}\label{def:b3:modules:companion}
Untuk $P = X^m + a_{m-1}X^{m-1} + \dots + a_0$ yang monik,
\emph{matriks pendamping}\index{matriks pendamping} adalah
\[
C_P = \begin{pmatrix}
0 & & & -a_0\\
1 & \ddots & & -a_1\\
& \ddots & 0 & \vdots\\
& & 1 & -a_{m-1}
\end{pmatrix} :
\]
yaitu matriks ``perkalian dengan $X$'' pada $K[X]/(P)$ dalam
basis $1, \bar X, \dots, \bar X^{m-1}$.
\end{definition}

\begin{theorem}[Frobenius: bentuk kanonik rasional]
\label{thm:b3:modules:frobenius}
Ada tepat satu barisan polinomial monik tak konstan $P_1
\mid P_2 \mid \dots \mid P_s$ (yaitu \emph{invarian
keserupaan}\index{invarian keserupaan} dari $u$) sedemikian sehingga, sebagai
modul-$K[X]$,
\[
V \cong K[X]/(P_1) \oplus \cdots \oplus K[X]/(P_s):
\]
dalam sebuah basis yang sesuai, $u$ mempunyai matriks diagonal blok
$\operatorname{diag}(C_{P_1}, \dots, C_{P_s})$. Lebih lanjut:
\begin{enumerate}
  \item $P_s = \mu_u$ (polinomial minimal) dan $P_1\cdots P_s =
        \chi_u$ (polinomial karakteristik); khususnya $\mu_u
        \mid \chi_u$ \emph{(Cayley--Hamilton yang dibuktikan ulang)} dan
        $\chi_u \mid \mu_u^{\,s}$, sehingga $\chi_u$ dan $\mu_u$ mempunyai
        faktor \omterm{def:b3:rings:divisibility}{tak tereduksi} yang sama.
  \item Dua endomorfisma (atau dua matriks persegi) serupa jika dan hanya jika
        \omterm{thm:b3:modules:frobenius}{invarian keserupaannya} sama.
\end{enumerate}
\end{theorem}

\begin{proof}
Teorema struktur (\cref{thm:b3:modules:structure}) yang diterapkan
pada \omterm{def:b3:rings:pidufd}{DIU} $K[X]$: \omterm{def:b3:modules:torsion}{modul torsi} $V$ terurai dengan
\omterm{thm:b3:modules:smith}{faktor invarian} $P_i$, yang dinormalkan menjadi monik (sebab unit $K[X]$ adalah
$K^\times$); tak ada bagian \omterm{def:b3:modules:free}{bebas} yang muncul (sebab $V$ bersifat \omterm{def:b3:modules:torsion}{torsi}). Pada setiap suku
siklik $K[X]/(P_i)$, perkalian dengan $X$ bermatriks $C_{P_i}$
dalam basis pangkat $\bar X$: menyambung basisnya memberikan
bentuk bloknya.

(1) Anihilator $V = \bigoplus K[X]/(P_i)$ adalah $(P_1)\cap
\dots\cap(P_s) = (P_s)$ (karena rantai keterbagiannya: $P_s$ merupakan kelipatan
persekutuan, dan kelas $1$ pada suku terakhir dimatikan
tepat oleh $(P_s)$): jadi $\mu_u = P_s$. Untuk $\chi_u$: pada suku
siklik berlaku $\chi_{C_P} = P$, lewat induksi pada $m = \deg P$. Menjabarkan
$\det(XI_m - C_P)$ sepanjang baris pertama (yang entrinya $X$, lalu
nol, lalu $a_0$ di kolom terakhir):
\[
\det(XI_m - C_P) = X\,\det\bigl(XI_{m-1} - C_{\tilde P}\bigr)
+ (-1)^{1+m}\,a_0\,\det L ,
\]
dengan $\tilde P = X^{m-1} + a_{m-1}X^{m-2} + \dots + a_1$ (berbentuk
sama, satu ukuran lebih kecil) dan $L$ segitiga dengan diagonal
$(-1, \dots, -1)$, sehingga $\det L = (-1)^{m-1}$. Menurut induksi, suku
pertamanya $X\tilde P$ dan suku keduanya $a_0$: totalnya
$X\tilde P + a_0 = P$ (kasus dasar $m=1$: $\det(X + a_0) = P$).
Determinan berkalian atas bloknya:
$\chi_u = \prod P_i$. Cayley--Hamilton: $\chi_u \in (\mu_u)$
sebab $P_s \mid$ setiap\dots\ sebaliknya setiap $P_i \mid P_s$, sehingga
$\chi_u = \prod P_i$ membagi $P_s^{\,s} = \mu_u^s$; dan $\mu_u =
P_s$ membagi $\chi_u$ sebagai salah satu faktornya.

(2) Endomorfisma yang serupa berarti struktur \omterm{def:b3:modules:module}{modul} yang sekawan, sehingga
invariannya sama (menurut ketunggalan pada
\cref{thm:b3:modules:structure}); sebaliknya invarian yang sama
memberikan modul-$K[X]$ yang isomorfik, dan isomorfisma \omterm{def:b3:modules:module}{modul}
tidak lain adalah bijeksi linear yang menganyam kedua endomorfismanya: yakni
sebuah keserupaan.
\end{proof}

\begin{corollary}[Keserupaan tak terpengaruh perluasan lapangan]
\label{cor:b3:modules:descent}
Misalkan $K \subseteq L$ dua lapangan dan $M, N \in M_n(K)$. Jika $M$ dan
$N$ serupa atas $L$, maka keduanya serupa atas $K$.
\end{corollary}

\begin{proof}
\omterm{thm:b3:modules:frobenius}{Invarian keserupaan} $M$ dihitung lewat rumus minor Smith
(\cref{thm:b3:modules:smith}) yang diterapkan pada matriks penyajian
$XI_n - M$ atas $K[X]$ --- memang modul-$K[X]$
$V_M = K^n$ mempunyai penyajian $XI_n - M$: pemetaan $K[X]^n \to V_M$,
$(Q_i) \mapsto \sum Q_i(M)e_i$, surjektif dengan kernel yang dibangun oleh
kolom $XI_n - M$ (lewat pemeriksaan langsung: modulo kolom-kolom
itu, setiap unsur $K[X]^n$ tereduksi menjadi vektor konstan,
dan vektor konstan terpetakan secara bijektif; soal akhir pekan
menguraikannya). FPB polinomial tidak berubah di bawah perluasan
lapangan: jika $d$ merupakan FPB monik di $K[X]$ dari keluarga $(f_j)$,
maka Bézout memberikan $d = \sum u_jf_j$ dengan $u_j \in K[X]$, sehingga setiap
pembagi persekutuan $f_j$ \emph{di $L[X]$} membagi $d$; dan karena $d$
sendiri pembagi persekutuan, ia juga FPB di $L[X]$. Karena itu
\omterm{thm:b3:modules:smith}{faktor invarian} $XI_n - M$, yang berupa hasil bagi FPB minor
berturutan, sama atas $K$ maupun atas $L$: jadi $M, N$ mempunyai
\omterm{thm:b3:modules:frobenius}{invarian keserupaan} yang sama atas $L$ jika dan hanya jika demikian atas $K$; simpulkan dengan
\cref{thm:b3:modules:frobenius}(2).
\end{proof}

\begin{theorem}[Bentuk Jordan, diturunkan ulang]\label{thm:b3:modules:jordan}
Andaikan $\chi_u$ terpecah atas $K$ (misalnya $K = \C$). Dengan menerapkan
penguraian atas \omterm{thm:b3:modules:structure}{pembagi elementer} pada $V$ (lihat bukti
\cref{thm:b3:modules:structure}) alih-alih \omterm{thm:b3:modules:smith}{faktor invarian}:
\[
V \cong \bigoplus_{\lambda, j} K[X]\big/\bigl((X -
\lambda)^{k_{\lambda,j}}\bigr),
\]
dan dalam basis $\bigl(\overline{(X-\lambda)^{k-1}}, \dots,
\overline{(X - \lambda)}, \bar 1\bigr)$ pada setiap sukunya, $u$ bekerja
sebagai blok Jordan $J_k(\lambda)$: jadi setiap endomorfisma dengan polinomial
karakteristik yang terpecah mempunyai basis Jordan, dan multihimpunan
bloknya $(\lambda, k)$ bersifat tunggal.\index{bentuk Jordan}
\end{theorem}

\begin{proof}
\omterm{thm:b3:modules:structure}{Pembagi elementer} \omterm{def:b3:modules:torsion}{modul torsi} $V$ adalah
$(X - \lambda)^k$ dengan $X - \lambda$ menjelajahi faktor \omterm{def:b3:rings:divisibility}{tak
tereduksi} $\mu_u$ (yang terpecah, sebab $\chi_u$ terpecah dan keduanya
mempunyai faktor \omterm{def:b3:rings:divisibility}{tak tereduksi} yang sama,
\cref{thm:b3:modules:frobenius}). Di $W = K[X]/((X-\lambda)^k)$,
tetapkan $f_j = \overline{(X - \lambda)^{k-j}}$ untuk $j = 1, \dots, k$:
maka $(X - \lambda)f_j = f_{j-1}$ (dengan $f_0 = 0$), yakni $u(f_j)
= \lambda f_j + f_{j-1}$: jadi matriks $u$ pada $(f_1, \dots,
f_k)$ tepat berupa $J_k(\lambda)$ (dengan angka satu di atas diagonalnya).
Ketunggalan multihimpunan \omterm{thm:b3:modules:structure}{pembagi elementernya} adalah
\cref{thm:b3:modules:structure}.
\end{proof}

\begin{remark}\label{rem:b3:modules:overview}
Hierarki bentuk kanoniknya kini menjadi jelas:
bentuk \emph{rasional} ada atas setiap lapangan dan mendeteksi
keserupaan secara mutlak (\cref{cor:b3:modules:descent}); sedangkan
bentuk \emph{Jordan} merupakan penghalusannya bila $\chi_u$ terpecah.
Bukti teorema Jordan pada Tahun ke-2 yang bersandar pada pencacahan dimensi
kini tercakup: seluruh kombinatorikanya ternyata aritmetika \omterm{def:b3:rings:pidufd}{DIU}
$K[X]$.
\end{remark}

\section{Latihan}

\begin{exercise}[$\star$]\label{exo:b3:modules:1}
(a) Tunjukkan bahwa $\Q$ tidak \omterm{def:b3:modules:free}{dibangun secara berhingga} sebagai modul-$\Z$.
(b) Tunjukkan bahwa $\Q$ \omterm{def:b3:modules:free}{bebas} \omterm{def:b3:modules:torsion}{torsi} tetapi tidak \omterm{def:b3:modules:free}{bebas}.
(c) Mengapa kedua pernyataan itu tidak bertentangan dengan
\cref{thm:b3:modules:structure}?
\end{exercise}

\begin{exercise}[$\star$]\label{exo:b3:modules:2}
Daftarkan semua grup abelian berorde $360$ sampai isomorfisma, baik dalam
bentuk \omterm{thm:b3:modules:structure}{pembagi elementer} maupun \omterm{thm:b3:modules:smith}{faktor invarian}. Ada berapa grup abelian
berorde $p^5$?
\end{exercise}

\begin{exercise}[$\star$]\label{exo:b3:modules:3}
Hitunglah \omterm{thm:b3:modules:smith}{bentuk normal Smith} atas $\Z$ dari
\[
B = \begin{pmatrix} 2 & 4\\ 6 & 8 \end{pmatrix},
\qquad
C = \begin{pmatrix} 2 & 0 & 0\\ 0 & 3 & 0\\ 0 & 0 & 12
\end{pmatrix},
\]
lalu kenali grup abelian $\Z^2/B\Z^2$ dan $\Z^3/C\Z^3$.
\end{exercise}

\begin{exercise}[$\star\star$]\label{exo:b3:modules:4}
Misalkan $L \subseteq \Z^n$ sebuah subgrup berrank $n$ yang basisnya berupa
kolom $B \in M_n(\Z)$, $\det B \neq 0$. Tunjukkan bahwa
$\Z^n/L$ berhingga dengan kardinalitas $\abs{\det B}$, dan bahwa $\Z^n/L
\cong \prod_i \Z/d_i\Z$ untuk \omterm{thm:b3:modules:smith}{faktor invarian} $d_i$ dari $B$.
Berilah gambaran dengan $L = \Z(2,0) + \Z(1,3)$.
\end{exercise}

\begin{exercise}[$\star\star$]\label{exo:b3:modules:5}
Misalkan $A$ daerah integral. (a) Periksalah bahwa $T(M)$ submodul dan
bahwa $M/T(M)$ \omterm{def:b3:modules:free}{bebas} \omterm{def:b3:modules:torsion}{torsi}. (b) Tunjukkan bahwa \omterm{def:b3:rings:ideal}{ideal} $(X, Y)$
dari $K[X,Y]$, sebagai modul-$K[X,Y]$, \omterm{def:b3:modules:free}{bebas} \omterm{def:b3:modules:torsion}{torsi} tetapi tidak \omterm{def:b3:modules:free}{bebas}:
jadi teorema struktur memang sungguh memerlukan hipotesis \omterm{def:b3:rings:pidufd}{DIU}.
\end{exercise}

\begin{exercise}[$\star\star$]\label{exo:b3:modules:6}
(a) Tunjukkan bahwa $2\Z$ tak punya komplemen langsung di dalam modul-$\Z$
$\Z$: submodul dari \omterm{def:b3:modules:free}{modul bebas} memang \omterm{def:b3:modules:free}{bebas}
(\cref{thm:b3:modules:submodule}), tetapi belum tentu menjadi suku
langsung. (b) Tunjukkan bahwa jika $M \subseteq A^n$ (dengan $A$ sebuah \omterm{def:b3:rings:pidufd}{DIU}) memenuhi:
$A^n/M$ \omterm{def:b3:modules:free}{bebas} \omterm{def:b3:modules:torsion}{torsi}, maka $M$ \emph{memang} suku langsung.
\end{exercise}

\begin{exercise}[$\star\star$]\label{exo:b3:modules:7}
Selesaikan di $\Z^2$ sistem
\[
\begin{cases}
2x + 4y \equiv b_1 \pmod{20},\\
6x + 8y \equiv b_2 \pmod{20},
\end{cases}
\]
lalu tentukan untuk pasangan $(b_1, b_2)$ yang mana penyelesaian itu ada, dengan memakai bentuk
Smith pada \cref{exo:b3:modules:3} (penggantian variabel yang terbalikkan
di kedua ruas).
\end{exercise}

\begin{exercise}[$\star\star$]\label{exo:b3:modules:8}
(a) Tentukan semua \omterm{thm:b3:modules:frobenius}{invarian keserupaan} dan \omterm{thm:b3:modules:jordan}{bentuk Jordan} yang mungkin
bagi matriks nilpoten $4 \times 4$, dipilah menurut partisi
$4$ yang diwujudkannya.
(b) Tampilkan dua matriks kompleks $4\times4$ dengan polinomial
karakteristik \emph{dan} polinomial minimal yang sama tetapi tidak
serupa, lalu buktikan bahwa untuk $n \leq 3$ hal itu tak mungkin terjadi.
\end{exercise}

\begin{exercise}[$\star\star\star$]\label{exo:b3:modules:9}
Misalkan $u \in \mathcal L(V)$, $\dim V = n$. Tunjukkan bahwa berikut ini
saling ekuivalen: (i) $V$ merupakan modul-$K[X]$ siklik (ada $x$
dengan $V = K[u]x$, yaitu sebuah \emph{vektor siklik}); (ii) $\mu_u = \chi_u$;
(iii) $s = 1$ pada \cref{thm:b3:modules:frobenius}. Simpulkan bahwa
\omterm{def:b3:modules:companion}{matriks pendamping} mempunyai vektor siklik, lalu tentukan kapan
matriks diagonal mempunyainya.
\end{exercise}

\begin{exercise}[$\star\star\star$]\label{exo:b3:modules:10}
Untuk $M \in M_n(\Z)$ yang dipandang sebagai endomorfisma $\Z^n$, buktikan
rumus indeksnya: jika $\det M \ne 0$, maka
$[\Z^n : M\Z^n] = \abs{\det M}$, lalu simpulkan bahwa $M \in GL_n(\Z)$
jika dan hanya jika $\det M = \pm 1$. Terapannya: grup $\Z^2$ mempunyai tepat
$\sigma_1(m) = \sum_{d \mid m} d$ subgrup berindeks $m$.
\emph{(Cacahlah matriks dalam bentuk Hermite $\bigl(\begin{smallmatrix}
a & b\\ 0 & d\end{smallmatrix}\bigr)$, $ad = m$, $0 \leq b <
d$.)}
\end{exercise}

\begin{exercise}[$\star\star$]\label{exo:b3:modules:11}
(Persamaan $x^k = e$ di dalam grup abelian) Misalkan $G$ grup abelian
berhingga dengan \omterm{thm:b3:modules:smith}{faktor invarian} $d_1 \mid d_2 \mid \dots
\mid d_s$.
(a) Tunjukkan bahwa untuk setiap $k \geq 1$,
\[
\#\{x \in G : x^k = e\} \;=\; \prod_{i=1}^{s}\gcd(k, d_i) .
\]
(b) Simpulkan: grup abelian berhingga bersifat siklik jika dan hanya jika
untuk setiap $k$, persamaan $x^k = e$ mempunyai paling banyak $k$
penyelesaian.
(c) Peroleh kembali kesiklikan subgrup berhingga dari $K^\times$
(dengan $K$ sebuah lapangan, \cref{ch:b3:galois}): mengapa polinomial
$X^k - 1$ menjamin kriteria pada (b)?
\end{exercise}

\begin{exercise}[$\star\star\star$]\label{exo:b3:modules:12}
(Matriks elementer membangun) (a) Tunjukkan bahwa $M \in M_n(\Z)$
terbalikkan di $M_n(\Z)$ jika dan hanya jika $\det M = \pm1$.
(b) Tunjukkan bahwa $SL_2(\Z)$ dibangun oleh dua matriks
elementer
$E = \bigl(\begin{smallmatrix}1 & 1\\ 0 &
1\end{smallmatrix}\bigr)$ dan
$F = \bigl(\begin{smallmatrix}1 & 0\\ 1 &
1\end{smallmatrix}\bigr)$.
\emph{(Jalankan algoritma Euclid pada kolom pertama $M
\in SL_2(\Z)$ lewat perkalian kiri dengan pangkat $E, F$,
hingga mencapai $\pm\bigl(\begin{smallmatrix}1 & *\\ 0 &
1\end{smallmatrix}\bigr)$; selesaikan sisanya dengan tangan --- perhatikan
bahwa $-I = (EF^{-1}E)^2$.)}
(c) Jelaskan kaitannya dengan reduksi Smith: atas $\Z$,
operasi baris dan kolom berdeterminan $1$ sudah cukup untuk
mendiagonalkan, sampai pada tandanya.
\end{exercise}

\section{Soal: komutan dan komutan ganda}

\begin{problem}[{Soal akhir pekan --- bentuk rasional, komutan,
bikomutan}]\label{pb:b3:modules:1}
Misalkan $K$ sebuah lapangan, $V$ ruang vektor-$K$ berdimensi $n \geq
1$, dan $u \in \mathcal L(V)$. Kita mempelajari \emph{komutan}
\[
\mathcal C(u) = \{v \in \mathcal L(V) : uv = vu\},
\]
yaitu subaljabar $\mathcal L(V)$ yang memuat $K[u] = \{P(u) : P \in
K[X]\}$, lalu kita buktikan rumus dimensi Frobenius dan teorema komutan
ganda: $\mathcal C(\mathcal C(u)) = K[u]$. Sepanjang soal ini,
$V$ adalah modul-$K[X]$ yang didefinisikan oleh $u$, dengan \omterm{thm:b3:modules:smith}{faktor invarian}
$P_1 \mid \cdots \mid P_s$ dan penguraian siklik $V =
\bigoplus_{i=1}^s V_i$, $V_i = K[u]\,x_i \cong K[X]/(P_i)$, $n_i
= \deg P_i$ (\cref{thm:b3:modules:frobenius}).

\medskip
\noindent\textbf{Bagian I --- Matriks penyajian $XI - M$, dan
pemanasan.}
\begin{enumerate}
  \item Misalkan $M \in M_n(K)$ dan misalkan $\varphi \colon K[X]^n \to V_M
        = K^n$ mengirim $(Q_1, \dots, Q_n)$ ke $\sum_i Q_i(M)e_i$.
        Tunjukkan bahwa $\varphi$ merupakan morfisma surjektif antar
        modul-$K[X]$ dan bahwa setiap kolom $XI_n - M$ terletak
        di $\ker\varphi$.
  \item Tunjukkan bahwa, modulo kolom $XI_n - M$, setiap
        unsur $K[X]^n$ kongruen dengan sebuah vektor konstan
        \emph{(turunkan derajatnya dengan memakai $Xe_i \equiv \sum_j m_{ji}e_j$)},
        lalu simpulkan $\ker\varphi = (XI_n - M)\,K[X]^n$: jadi \omterm{def:b3:modules:module}{modul}
        $V_M$ bermatriks penyajian $XI_n - M$. Peroleh kembali
        titik tolak \cref{cor:b3:modules:descent}: bahwa
        \omterm{thm:b3:modules:frobenius}{invarian keserupaan} $M$ adalah \omterm{thm:b3:modules:smith}{faktor invarian}
        $XI_n - M$ yang bukan unit.
  \item Hitunglah \omterm{thm:b3:modules:frobenius}{invarian keserupaan} dari: matriks skalar
        $\lambda I_n$; matriks diagonal dengan entri diagonal
        yang berbeda-beda; blok Jordan $n \times n$ $J_n(0)$;
        dan $\operatorname{diag}(J_2(0), J_1(0))$ untuk $n = 3$.
  \item Tunjukkan bahwa $\dim K[u] = \deg \mu_u = n_s$.
\end{enumerate}

\medskip
\noindent\textbf{Bagian II --- Morfisma antar \omterm{def:b3:modules:module}{modul} siklik.}
\begin{enumerate}[resume]
  \item Misalkan $P, Q$ monik dan tak konstan. Tunjukkan bahwa sebuah
        morfisma-$K[X]$ $f \colon K[X]/(P) \to K[X]/(Q)$
        ditentukan oleh $f(\bar 1)$, dan bahwa $\bar c \in K[X]/(Q)$
        dapat berperan sebagai $f(\bar 1)$ jika dan hanya jika $P\bar c = 0$ di $K[X]/(Q)$.
  \item Simpulkan
        \[
        \operatorname{Hom}_{K[X]}\bigl(K[X]/(P),\, K[X]/(Q)\bigr)
        \;\cong\; K[X]\big/\bigl(\gcd(P, Q)\bigr),
        \]
        yang berdimensi $\deg \gcd(P, Q)$ atas $K$. \emph{(Tunjukkan bahwa
        penyelesaian $\bar c$ dari $P\bar c = 0$ di $K[X]/(Q)$
        membentuk submodul siklik yang dibangun oleh
        $\overline{Q/\gcd(P,Q)}$.)}
  \item Buktikan \emph{rumus Frobenius}:
        \[
        \dim_K \mathcal C(u) = \sum_{i,j=1}^{s} \deg\gcd(P_i,
        P_j) = \sum_{i=1}^{s} (2s - 2i + 1)\, n_i .
        \]
        \emph{(Sebuah $v$ yang komutatif tidak lain adalah endomorfisma-$K[X]$ dari
        $V$; uraikan $\operatorname{End}(\bigoplus_i V_i)$ sebagai
        matriks morfisma $V_j \to V_i$ lalu pakailah
        rantai keterbagiannya.)}
  \item Simpulkan $\dim \mathcal C(u) \geq n$, dengan kesamaan jika dan hanya jika $u$
        siklik ($s = 1$), lalu hitunglah $\dim\mathcal C(u)$ untuk
        $u = \lambda\,\mathrm{id}$: keduanya ujung ekstrem rumus itu.
  \item Periksalah rumus Frobenius secara langsung untuk
        $\operatorname{diag}(J_2(0), J_1(0))$ dengan menghitung
        komutannya secara eksplisit sebagai matriks $3\times3$.
\end{enumerate}

\medskip
\noindent\textbf{Bagian III --- Teorema komutan ganda.}
Misalkan $w \in \mathcal C(\mathcal C(u))$; kita buktikan $w \in K[u]$.
\begin{enumerate}[resume]
  \item Tunjukkan $K[u] \subseteq \mathcal C(\mathcal C(u))$, dan bahwa
        setiap $w \in \mathcal C(\mathcal C(u))$ komutatif dengan $u$
        --- jadi inklusi yang hendak dibuktikan,
        $\mathcal C(\mathcal C(u)) \subseteq K[u]$, sungguh
        mempertajam $w \in \mathcal C(u)$.
  \item Andaikan lebih dahulu bahwa $u$ bersifat \emph{siklik}, $V = K[u]x$.
        Tunjukkan secara langsung bahwa $\mathcal C(u) = K[u]$
        \emph{(masukkan sebuah $v$ yang komutatif pada $x$: $v(x) = P(u)x$
        untuk suatu $P$, lalu bandingkan $v$ dengan $P(u)$ pada basis
        $u^k x$)}, lalu simpulkan teoremanya untuk kasus ini.
  \item Kembali ke kasus umumnya. Untuk setiap $i$, misalkan $\pi_i\colon
        V \to V_i$ proyeksi sepanjang suku lainnya.
        Tunjukkan $\pi_i \in \mathcal C(u)$, lalu simpulkan bahwa $w$
        mempertahankan setiap $V_i$ dan komutatif dengan $u_i =
        u\restriction_{V_i}$; simpulkan lewat pertanyaan 11 yang diterapkan pada
        $u_i$ yang siklik: ada polinomial $Q_i$ dengan
        $w\restriction_{V_i} = Q_i(u)\restriction_{V_i}$.
  \item Tersisa merekatkan $Q_i$ menjadi satu polinomial. Untuk $i
        \leq j$ (sehingga $P_i \mid P_j$), tunjukkan bahwa $\eta_{ij} \colon
        V_j \to V_i$, $R(u)x_j \mapsto R(u)x_i$, merupakan
        morfisma-$K[X]$ yang terdefinisi dengan baik \emph{(yang perlu diperiksa
        adalah bahwa $R(u)x_j = 0$ mengakibatkan $R(u)x_i = 0$)}, dan bahwa
        $\tilde\eta_{ij} = \eta_{ij}\circ\pi_j$, yang diperluas dengan $0$
        pada suku lainnya, terletak di $\mathcal C(u)$.
  \item Dengan memakai $w\tilde\eta_{ij} = \tilde\eta_{ij}w$, tunjukkan $Q_i
        \equiv Q_j \pmod{P_i}$ untuk $i \leq j$. Simpulkan bahwa $Q =
        Q_s$
        memenuhi $Q \equiv Q_i \pmod {P_i}$ untuk setiap $i$, sehingga
        $w = Q(u)$ pada setiap $V_i$, dan karenanya pada $V$:
        \[
        \boxed{\ \mathcal C(\mathcal C(u)) = K[u].\ }
        \]
  \item (Penutup) Simpulkan dari teorema itu: jika $v$ komutatif dengan
        \emph{setiap} matriks yang komutatif dengan $u$, dan $u$
        siklik, maka $v$ merupakan polinomial dalam $u$; lalu berilah
        contoh yang menunjukkan bahwa $\mathcal C(u) = K[u]$ gagal untuk $u =
        \mathrm{id}$, $n \geq 2$ --- di manakah tepatnya
        kesiklikan itu masuk?
\end{enumerate}

\medskip
\noindent\textbf{Bagian IV --- Hasil panen \omterm{thm:b3:modules:frobenius}{invarian
keserupaan}.} Bentuk kanonik rasional adalah sebuah mesin; berikut
lima keluarannya yang klasik.
\begin{enumerate}[resume]
  \item (Transpos) Tunjukkan bahwa setiap $M \in M_n(K)$ serupa dengan
        transposnya ${}^tM$. \emph{(Operasi yang membawa
        $XI - M$ ke bentuk Smith, setelah ditranspos, membawa $XI - {}^tM$ ke
        bentuk Smith yang sama: jadi \omterm{thm:b3:modules:frobenius}{invarian keserupaannya} sama.)}
  \item (Penurunan keserupaan) Misalkan $K \subseteq L$ sebuah perluasan
        lapangan dan $M, N \in M_n(K)$. Tunjukkan bahwa jika $M$ dan $N$
        serupa atas $L$, maka keduanya serupa atas $K$.
        \emph{(Bentuk Smith $XI - M$ yang dihitung di $K[X]$
        tetap merupakan bentuk Smith di $L[X]$ --- mengapa \omterm{thm:b3:modules:smith}{faktor invariannya}
        tidak berubah?)} Akibat yang layak dihafal: dua matriks
        real yang sekawan di $GL_n(\C)$ juga sekawan di
        $GL_n(\R)$.
  \item (Klasifikasi nilpoten) Misalkan $u$ nilpoten. Tunjukkan
        bahwa banyaknya blok berukuran $\geq k$ pada
        penguraiannya atas blok Jordan nilpoten sama dengan
        $\operatorname{rk}u^{k-1} - \operatorname{rk}u^k$, lalu
        simpulkan: kelas nilpoten di $M_n(K)$, untuk lapangan $K$ apa pun,
        berkorespondensi satu-satu dengan partisi $n$. Ada berapa
        kelas nilpoten di $M_5(K)$?
  \item (Sepasang contoh konkret) Tentukan \omterm{thm:b3:modules:frobenius}{invarian keserupaan}
        operator turunan $D\colon P \mapsto P'$ yang bekerja pada ruang
        $K_{n-1}[X]$ berisi polinomial berderajat $< n$: (a) untuk $K =
        \Q$; (b) untuk $K = \mathbb F_p$ dengan $p < n$
        \emph{(pada karakteristik $p$ berlaku $(X^p)' = 0$: hitunglah
        $\ker D^k$ lalu pakailah pertanyaan 18)}.
  \item (\omterm{ex:b3:groups:actions}{Kelas konjugasi} $GL_2(\mathbb F_q)$) Dengan memakai
        \omterm{thm:b3:modules:smith}{faktor invarian}, tunjukkan bahwa setiap kelas
        $GL_2(\mathbb F_q)$ tepat berjenis salah satu dari empat berikut:
        pusat $aI$; dapat didiagonalkan dengan dua nilai eigen
        berbeda $a \neq b$ di $\mathbb F_q^\times$;
        tak semisederhana dengan polinomial minimal $(X - a)^2$;
        dan siklik dengan polinomial karakteristik yang \omterm{def:b3:rings:divisibility}{tak tereduksi}.
  \item Cacahlah kelas pada setiap jenisnya lalu simpulkan:
        $GL_2(\mathbb F_q)$ mempunyai tepat $q^2 - 1$ \omterm{ex:b3:groups:actions}{kelas
        konjugasi}. \emph{(Cacahlah polinomial kuadrat monik yang \omterm{def:b3:rings:divisibility}{tak tereduksi} atas
        $\mathbb F_q$; lalu pasangan tak terurut $\{a, b\}$; ingatlah bahwa
        keterbalikan membatasi suku konstannya.)}
  \item (Siklik itu generik) Tunjukkan bahwa $M \in M_2(\mathbb F_q)$
        gagal menjadi siklik jika dan hanya jika $M$ skalar, lalu simpulkan bahwa
        matriks $2\times2$ acak seragam atas $\mathbb F_q$ bersifat
        siklik dengan peluang $1 - q^{-3}$. Nyatakan heuristik
        serupa untuk $M_n$ dan $q$ besar (tanpa perlu bukti):
        matriks tak siklik itu langka --- itulah sebabnya
        Bagian III \cref{pb:b3:modules:1} baru memerlukan kerja sungguhan
        setelah melewati kasus generiknya.
\end{enumerate}

\medskip
\noindent\textbf{Bagian V --- Pelengkap.}
\begin{enumerate}[resume]
  \item (Pusat komutan) Tunjukkan bahwa pusat aljabar
        $\mathcal C(u)$ tepat sama dengan $K[u]$ \emph{(gabungkan
        kedua inklusi pada Bagian III)}. Simpulkan bahwa
        $\mathcal C(u)$ komutatif jika dan hanya jika $u$ siklik ---
        sehingga kasus kesamaan pada pertanyaan 8 diperoleh kembali lewat jalur
        yang murni struktural.
  \item (Dimensi mana yang muncul?) Simpulkan dari rumus
        Frobenius bahwa $\dim\mathcal C(u) \equiv n \pmod 2$ untuk
        setiap $u$. Lalu tentukan himpunan nilai persis yang diambil
        oleh $\dim \mathcal C(u)$ ketika $u$ menjelajahi
        $\mathcal L(V)$ dengan $\dim V = 4$: tunjukkan bahwa himpunan itu $\{4, 6,
        8, 10, 16\}$ \emph{(cacahlah barisan derajat $n_1
        \leq \dots \leq n_s$ yang berjumlah $4$ lalu wujudkan masing-masing lewat
        sebuah matriks nilpoten)}. Khususnya $12$ dan $14$, meski
        berkeparitasan benar, tidak tercapai: jadi syarat keparitasan
        itu perlu tetapi tidak cukup.
  \item (\omterm{cor:b3:groups:classeq}{Persamaan kelas} $GL_2(\mathbb F_3)$) Untuk $q = 3$,
        hitunglah ukuran setiap \omterm{ex:b3:groups:actions}{kelas konjugasi} pada pertanyaan 20
        lewat orbit--stabilisator: \omterm{ex:b3:groups:actions}{sentralisator} sebuah $M$ yang
        \emph{siklik} di $GL_2(\mathbb F_q)$ adalah grup unit
        dari $K[M]$ (pertanyaan 11). Kenali $K[M]$ pada
        ketiga jenis yang bukan pusat, daftarkan ketiga polinomial kuadrat
        monik yang \omterm{def:b3:rings:divisibility}{tak tereduksi} atas $\mathbb F_3$, lalu periksalah persamaan
        kelasnya
        \[
        48 = \abs{GL_2(\mathbb F_3)}
        = 2\cdot1 + 1\cdot12 + 2\cdot8 + 3\cdot6,
        \]
        dengan $2 + 1 + 2 + 3 = 8 = q^2 - 1$ kelas, seperti
        diramalkan pertanyaan 21.

\end{enumerate}
\end{problem}
